\documentclass{report}

\usepackage{fancyhdr} % Cabeceras de página
\usepackage{lastpage} % Módulo para añadir una referencia a la última página
\usepackage{titling} % No tengo claro para qué es esto
\usepackage[left=2cm,right=2cm,top=3cm,bottom=2cm]{geometry} % Márgenes
\usepackage[T1]{fontenc}
\usepackage[utf8x]{inputenc}
\usepackage{xspace}
\usepackage{graphicx}
\usepackage{tikz}
\usepackage{float}

\setcounter{secnumdepth}{3}


\title{Fault Manager Lite (Technical report)}
\date{\today}
\author{SGNIP company}

\fancyhf{}
\fancypagestyle{plain}{%
	\lhead{\small \itshape \thetitle\, -\, \thedate\, -\, PINGS}
	\rhead{\vspace{-20pt} \includegraphics[width =40 pt]{../Logo.jpg}}
	\cfoot{\thepage\ of \pageref{LastPage}}
	\rfoot{}
}

\begin{document}
\maketitle
\tableofcontents
\newpage
\chapter{Abstract}

\chapter{Introduction}

The Autonomous University of Madrid (UAM) has reported the numerous problems it has on detecting faults that arise on its campus and its facilities, whose reparation usually takes excessive time and is poorly organised. A late detection of faults in the facilities delays its reparation, stopping its users to continue using them as normal and complicating the maintenance staff labour. Regarding the wishes of the campus users and realising the maintenance problems it has, the UAM has organised a contest to choose the best project proposal that solves them.
This is where our organization, Triforce, enters the scene: we have analyzed the problem exhaustively and designed a web application, Fault Manager Lite (FML), that meets all the requirements expected, solves the problems and also includes new extra features that makes it even more useful.

\paragraph{Purpose: }
The purpose of this Software Requirements Specification (SRS) document is to provide a
detailed description of the functionalities of the FML system. This document will cover each of the system's intended features, as well as offer a preliminary glimpse of the software application's User Interface (UI).

\section{Structure}

\section{Definitions and abbreviations}

\begin{itemize}
\item \textbf{SRS: } Software Requirements Specification.

\item \textbf{FML: } Fault Manager Lite. This is the System's name.

\end{itemize}


\section{Methodology}

The methodological procedure followed for preparing the SRS of the software system FML we have developed includes the following techniques:
\begin{itemize}
\item[1] Deep analysis of the information given by the UAM: the potential maintenance problem it has, the causes that lead to this problem and various aspects we must consider before proceeding to the next step (web application;  works on PCs, tablets and smartphones; target users; basic purpose).
\item[2] Application of the technique of brainstorming to generate general ideas for the application. Brainstorming was then applied to these general ideas in order to give shape to them. The results of the brainstorm carried out are cleaned up in the Annex A. This step resulted on the detection of 5 subsystems.
\item[3] Research on the Internet about other similar systems, analyzing the functionality of these systems and describing them in a structured way by specifying their advantages and disadvantages, and extracting good ideas to incorporate to our own project, increasing its market value. These ideas also helped putting the final touches to our ideas from the brainstorm.
\item[4] Interview with one member of the UAM's technical staff. In this interview we clarified the obscure ideas we still had, and modified some of our previous ideas to adapt them to the answers given. The answers to the interview were well-considered; as the technical staff will be the main user of our application, their point of view is important to us.
\item[5] Departing from the ideas obtained, defined the requirements (functional and non-functional) of the application and design attractive mock-ups that fulfill them. These mock-ups are only a demo and may change in the final version of the application.
\end{itemize}

\chapter{Project Definition}

FML is based on the jolly cooperation between the users of the facilities of the UAM campus and the maintenance staff in charge of them.

As the users will be the ones that will detect the sooner faults on the facilities the campus offers them, they are also the most suitable to report the problems they are having, in order to getting them fixed as soon as possible. With the FML web application, it will only take less than a minute to fill the form and send it to the maintenance staff.

Using the reports of faults detected in the campus, the maintenance personnel will stop losing its precious time revising the installations looking for faults and will be able to focus on the repairs. Apart from this benefit, the maintenance staff will also have a better way to coordinate efforts, as the FML will provide an automatic assignment system to assign repairs to each of the members avoiding overloading any of them and taking into account their distance to the problem, saving time on displacements.

In summary, the users will be able to easily report faults on facilities they are using in order to have them fixed in the less time possible, while the maintenance personnel will multiply its current performance as the majority of their resources will stop being wasted on revisions, but on repairs. We state that the FML web application is the answer to the UAM problems, as we will show you with this SRS.


\section{Goals and functionalities}


Our main goal is the development of an extreme-functional, bug-free web application, that runs either on PCs, tablets or smartphones and facilitates the labor of the maintenance staff of the UAM.

Based on the jolly cooperation between the members of the UAM community (students, teaching and research staff and administration and services personnel) and the maintenance personnel of the facilities and installations of the UAM campus, this application has been designed to solved the potential maintenance problems that the UAM has declared to have.

The problems we aim to alleviate with this application include the following:


\begin{itemize}
\item Difficulty of detecting faults on facilities of the campus.
\item Late detection of the problems, which implies late repairs of them.
\item Excessive resources and time wasted on revisions looking for potential faults.
\item Bad coordination between members of the technical staff.
\item Frequently overloading of some repairmen because of bad coordination.
\item Inability of a repairman to instantly inform that a fault has just been solved.
\item Difficulty of the users to report faults to the maintenance services.
The report system lacks of mobility functionality.
\item Lack of real-time visualization of pending and finished repairs.
\item Lack of real-time visualization of assigned tasks.
\item No faults history.
\item No statistics of faults.
\end{itemize}

\section{Subsystems or Modules}

The FML system has been designed to resolve these problems making use of a user-friendly interface, avoiding unnecessary or distracting buttons or effects. This system has been divided in several modules listed below, each of those offering one specific functionality and, in conjunction, solving the problems mentioned above:
\begin{itemize}
\item Task Manager:
\item Report System:
\item Notifications and Messaging System:
\item Users and Profile Managers:
\item Faults History.
\end{itemize}

There are several applications 'on the market' which are similar to our proposed app in some way or another and they could represent real and competitive alternatives to ours. However, none of these web applications can offer all the features FML offers (furthermore, we plan on adding some features which are not currently available in any of these applications), so our system is indeed the best solution proposed to solve the problems the UAM maintenance staff has to deal with. After our research of the Internet, a comparison of these competitors applications is gathered in Annex B.

\chapter{Catalog of Requirements}

\section{Document Service Requirements}

\subsection{Functional}

\subsubsection{General Requirements}

This requirements are defined for all roles inside FML (Admin, maintenance personnel and user)

\paragraph{Requirements focused on users}


\subparagraph{Log in: } Everyone must be able to login using UAM credentials. We will use the service provide by UAM to authenticate an user.

If the user is login in a smartphone, the system shouldn't ask for the password more than once (except for some specific operations defined in the role-depending requirements, user role category (\ref{Specifics_Secure_Requirements_for_user}))

\subparagraph{Lightweight application: } FML will be lightweight enough to be ran in Internet browser 4 years old.

\subparagraph{Faults definition: } Each fault will have 4 possible states: \textit{Pending to assign, assigned, solved, impossible}. Impossible state should be only categorized by an administrator.

\paragraph{Requirements focused on maintenance personnel}

\subparagraph{Categorizing faults: } There should be categories to tag each fault, depending on the difficulty, the material needed to be fixed.

\subparagraph{Fault's assignments: } FML will assign to maintenance personnel faults to be fixed (depending on the fault's category and maintenance person load of work and capability to solve the fault)

This will be set automatically and the system's administrators will be able to reassign the faults as they pleased.

\subparagraph{Task managing: } Everyone from the maintenance personnel will have read access to the faults data base. They will just have write access to faults they have been assigned to. Write access is needed so the maintenance person change fault's state to \textit{solved} or to \textit{in progress}.

\subparagraph{Communication: } Maintenance person assigned to fix a fault should be able to ask the fault's reporter for more information.

\subparagraph{When a fault is fixed: } The maintenance person assigned will mark as solved the fault. At that time, the reporter will get notify and thanked.

\subparagraph{Duplicated faults: } This is a big problem to be solved. What we propose to solve it has 2 aspects.
\begin{itemize}
\item FML will be able to mark as \textit{possible duplicated}, so all possible duplicates will be assign to the same maintenance people, so he can verify if they are duplicated or not.

\item On the other side, to prevent duplicated faults reports, FML will show a messages showing possible duplicates faults to the reporter. If the reporter marks the fault he is reporting as duplicated, he will earn some points too and will be added into the list of people to be notified when the task is solved.
\end{itemize}

\subsubsection{Role-depending requirements}

\paragraph{Reporter: } FML should provide an user the following abilities:
\begin{itemize}
\item Reporting a fault with the following information: \textit{Location, photograph (optional) and description}
\item See and check reported faults and see its state.
\item Answer questions the maintenance personnel will ask about the reported fault.
\item The specifics requirements where re-authentication is needed are: \textit{Answering questions asked by maintenance personnel}   \label{Specifics_Secure_Requirements_for_user}
\end{itemize}


\subsection{Non-Functional}

\subsubsection{General Requirements}

This requirements are defined for all roles inside FML (Admin, maintenance personnel and user)

\paragraph{Profile: } Everyone should have a profile with some private information such us name.
\subparagraph{Statistics: } In the profile page each user should be able to see the faults he has reported and it's state.

Plus, he should be able to see his conversations with maintenance personnel.


\subsubsection{Role-depending requirements}

\paragraph{User Role: }
FML should provide an user the following abilities:
\begin{itemize}
\item The interface to report a fault should be easy enough to let the user report the fault without losing much time.
\item The amount of points earned should be visible and the information of how much points does the user need to earn ETCS credit should be visible inside FML.
\end{itemize}


\chapter{Design of Mock-ups}

\chapter{Conclusions}


\appendix

\chapter{Brainstorming}

% -*- root: ../Technical_report.tex -*-
\section{Ideas about new development}

\subsection{Roles of the Application}

We define the user as the reporter of the incidence and the maintenance personnel as the people in charge of fixing the faults.

Every one uses the same application, but it will be shown a different menu depending on the role after the authentication.

\begin{itemize}
\item \textbf{Admin role: } The person or people in charge of banning and modifying manually what was automatically set by the system.

\item\textbf{Maintenance person:} You can see as a Maintenance person what task has been assigned to you and you can mark them as finished or as a false fault.

You can also ask for more information about the fault if it's not precise enough.

\item \textbf{User role: }  We may define subroles such as professor, student or workers (waiters, secretaries ... ) because of the liability of each group of people.

\end{itemize}

Everyone is a user and everyone can report faults but not  every one can be Maintenance person or user.

\subsection{Task priority}
We talked about 2 options:

\begin{itemize}
\item The user sets the priority of the task and the system has to take on account the liability of the user to set a real priority to the task.

\item There are tags and categories the user can add to the report and the system computes a real priority to the task.

Within this idea we need to add \textit{others} category, just in case.
\end{itemize}

\subsection{Location}

We have to distinguish 2 cases, if you are inside a building or not.

If you are not inside a building, the location will be automatic by GPS technology, but when faults are relative to inside of buildings material it will be necessary to ask the user the exact location. This will be ask by an interactive map.


We will also use WiFi location because sometimes GPS is a slow technology to be set and reliable.

\subsection{Manage task assignation}
The system automatically assigns tasks, depending on the location of the faults and the Maintenance person and the availability of the maintenance personnel.

Admin can set manually the task assignation.

We should take care of the following situation:

A Maintenance person works a lot and fix every task he was assigned, and, because he is free, the system assigns some more. On the other side, a Maintenance person is not working. As he has lot of pending task, system will not add more task to him. 

Plus, we can't just count the number of faults fixed, because there are ones more difficult than others. May be we should add \textit{difficulty of fault} depending on category to try control the how much job you have been assigned.

\subsection{Graphical interface}
We have Guille's application in which we will add just a profile page in which users can see
\begin{itemize}
\item the false faults they have reported .
\item the fault's reported state.
\end{itemize}

Every user's profile will be visible to the admin but he won't be able to see any personal data.

The admin can see profiles so before he can ban an user, he can check how many false and repeated faults he has reported.


\paragraph{Maintenance person interface}
The interface will show by color code the faults assignments classifying  they by pending, solved or not started.


\subsection{Repeated faults}
When an user is going to report a fault, the system will suggest if it's repeated showing possible faults. If the users says it's a new fault, it will be still marked as a possible repeated fault.

When a task is assigned to a Maintenance person, all possible repeated faults will be assigned to the same Maintenance person. It's Maintenance person's responsibility to say if they are the same or not.

If they are not the same, the task will be enqueued and the system will automatically be assigned to a Maintenance person, maybe the himself or may be not.


\subsection{Communication}

When a fault is fix, the system will notify the user who reported it.

\section{Questions to ask}

\begin{itemize}
\item When an emergency occurs in UAM (e.g., a fire, a flood...), does everyone get an alert?
\item Gamified experience? Do we ``reward'' the user for reporting faults or shall we instead offer detailed reports so the IT department can reward users as they want? (Credit's ETCS reward is possible, in the same way as languages)
\item We should take care of the following situation:

A Maintenance person works a lot and fix every task he was assigned, and, because he is free, the system assigns some more. On the other side, a Maintenance person is not working. As he has lot of pending task, system will not add more task to him. 

If he can assure us every one will work normally, that will solve the problem.
\end{itemize}




\chapter{Evaluate current technology}
\newpage

\chapter{Meetings}

% -*- root: ../Technical_report.tex -*-
\section{Meeting Announcements}
\subsection{Meeting Announcement - 1}

\textbf{From: } Iván Márquez Pardo

\textbf{To: } Guillermo Julián Moreno, Iván Márquez Pardo y Víctor de Juan Sanz.

\textbf{Date, Time and Place: } 27/01/2015 from 14:35 ultil we finish Brainstorming session in Science building, module 17.

\textbf{Purpose} Finish brainstorming session.

\subsection{Meeting Announcement - 2}


\textbf{From: }

\textbf{To: } Guillermo Julián Moreno, Iván Márquez Pardo y Víctor de Juan Sanz.


\textbf{Date, Time and Place: }


\textbf{Purpose: }


\textbf{Agenda: }

\textbf{Decision Follow-up: }


\textbf{Documentation: }


\section{Meetings Minutes}

We have excluded all in-class meetings. Should we add them?

\subsection{Meeting Minute - 1}


\textbf{Date and Time:} 27/01/2015 from 14:35 to 15:03.


\textbf{Participants: } Iván Márquez Pardo, Guillermo Julián Moreno y Víctor de Juan Sanz.


\textbf{Topics: }
Finish brainstorming and brainwriting session.


\textbf{Decitions Made: \\}

\begin{tabular}{|c|c|c|}
\hline Activity & Person Responsible & DeadLine \\\hline
Write down clearly all the ideas got from the brainstorming session & Víctor de Juan & 29/01/2015\\\hline

\end{tabular}

Brainstorm is finished, we have enough ideas to go to the next step.

\subsection{Meeting Minute - 2}


\textbf{Date and Time:}


\textbf{Participants: }

Iván Márquez Pardo, Guillermo Julián Moreno y Víctor de Juan Sanz.


\textbf{Topics: }

\textbf{Decitions Made: \\}

\begin{tabular}{|c|c|c|}
\hline Activity & Person Responsible & DeadLine \\\hline
olakase & olakase & olakase \\\hline

\end{tabular}
\section{Meeting Announcements}
\subsection{Meeting Announcement - 1}

\textbf{From: } Iván Márquez Pardo

\textbf{To: } Guillermo Julián Moreno, Iván Márquez Pardo y Víctor de Juan Sanz.

\textbf{Date, Time and Place: } 27/01/2015 from 14:35 ultil we finish Brainstorming session in Science building, module 17.

\textbf{Purpose} Finish brainstorming session.

\subsection{Meeting Announcement - 2}


\textbf{From: }

\textbf{To: } Guillermo Julián Moreno, Iván Márquez Pardo y Víctor de Juan Sanz.


\textbf{Date, Time and Place: }


\textbf{Purpose: }


\textbf{Agenda: }

\textbf{Decision Follow-up: }


\textbf{Documentation: }


\section{Meetings Minutes}

We have excluded all in-class meetings. Should we add them?

\subsection{Meeting Minute - 1}


\textbf{Date and Time:} 27/01/2015 from 14:35 to 15:03.


\textbf{Participants: } Iván Márquez Pardo, Guillermo Julián Moreno y Víctor de Juan Sanz.


\textbf{Topics: }
Finish brainstorming and brainwriting session.


\textbf{Decitions Made: }
Víctor de Juan has to write clearly all the ideas in a document.

Brainstorm is finished, we have enough ideas to go to the next step.

\subsection{Meeting Minute - 2}


\textbf{Date and Time:}


\textbf{Participants: }

Iván Márquez Pardo, Guillermo Julián Moreno y Víctor de Juan Sanz.


\textbf{Topics: }


\textbf{Decitions Made: }


\end{document}
